%======================================================================
\chapter{The homeomorphism from a zigzag}\label{ch:zigzag}
%======================================================================

The homeomorphism of Theorem~10 is the limit of a ``zigzag'': chains on the two spaces, refining each other alternately through chain maps. This chapter proves
that such a zigzag defines a homeomorphism (Bing and Jones, pp.~188--189), independently of how it was built. It is the analogue of \vtwo{Chapter~5}, where
two descriptions of pseudo-arcs gave a homeomorphism.

\section{The data}

Let $X$, $Y$ be metric spaces and $P\subseteq X$, $Q\subseteq Y$ compact. Suppose that for every $i\ge0$ we are given:
\begin{itemize}
\item a chain $D^i$ in $X$ covering $P$, every link meeting $P$, of mesh $<\varepsilon_i$;
\item a BJ-chain $E^i$ in $Y$ covering $Q$, every link meeting $Q$, of mesh $<\varepsilon_i$;
\item a chain map $H_i$ of $D^i$ onto $E^i$ and a chain map $K_i$ of $E^{i+1}$ into $D^i$;
\end{itemize}
where $\varepsilon_i\to0$, and the two conditions
\begin{description}
\item[\textnormal{(B)}] $D^{i+1}_j\subseteq D^i_{K_i(H_{i+1}(j))}$ for every link $D^{i+1}_j$;
\item[\textnormal{(B$'$)}] $E^{i+1}_j\subseteq E^i_{H_i(K_i(j))}$ for every link $E^{i+1}_j$.
\end{description}
\[
\begin{array}{ccccccc}
D^0 & \longleftarrow & D^1 & \longleftarrow & D^2 & \longleftarrow & \cdots\\
\downarrow{\scriptstyle H_0} & \nwarrow{\scriptstyle K_0} & \downarrow{\scriptstyle H_1} & \nwarrow{\scriptstyle K_1} & \downarrow{\scriptstyle H_2} & & \\
E^0 & \longleftarrow & E^1 & \longleftarrow & E^2 & \longleftarrow & \cdots
\end{array}
\]
For $p\in P$ let
\[
\st_i(p)=\bigcup\set{E^i_{H_i(j)}:\ p\in D^i_j},
\]
the union of the images of the links of $D^i$ containing $p$ (Bing and Jones write $H_id(p,i)$).

\formal{\lean{HPC.BJ59.starImage}.}

\section{The homeomorphism}

\begin{lemma}\label{lem:star-props}
For $p\in P$ and every $i$:
\begin{enumerate}[label=(\alph*)]
\item $\st_i(p)$ is one link of $E^i$ or the union of two adjacent ones; it is nonempty, meets $Q$, and has diameter $<2\varepsilon_i$;
\item $\st_{i+1}(p)\subseteq\st_i(p)$.
\end{enumerate}
\end{lemma}

\begin{proof}
(a) Any two links of $D^i$ containing $p$ meet, so they are adjacent; three links cannot be pairwise adjacent, so $p$ lies in one or two adjacent links. Their images
under the chain map $H_i$ are one link or two adjacent links. The links meet $Q$, and two meeting sets of diameter $<\varepsilon_i$ have a union of diameter $<2\varepsilon_i$.

(b) Let $p\in D^{i+1}_j$ and $k=K_i(H_{i+1}(j))$. By (B), $p\in D^i_k$, so $E^i_{H_i(k)}\subseteq\st_i(p)$. By (B$'$) applied to the link $E^{i+1}_{H_{i+1}(j)}$,
$E^{i+1}_{H_{i+1}(j)}\subseteq E^i_{H_i(K_i(H_{i+1}(j)))}=E^i_{H_i(k)}$.
\end{proof}

\formal{\lean{HPC.BJ59.mem_starImage}, \lean{HPC.BJ59.starImage_succ_subset}, \lean{HPC.BJ59.starImage_dist_le}.}

\begin{theorem}[Bing and Jones, pp.~188--189]\label{thm:zigzag}
Under the hypotheses above there is a homeomorphism $h\colon P\to Q$ characterized by: for $p\in P$ and $q\in Q$, $q=h(p)$ if and only if $q\in\cl{\st_i(p)}$ for every $i$.
\end{theorem}

\begin{proof}
\emph{Definition.} By Lemma~\ref{lem:star-props} the sets $\cl{\st_i(p)}\cap Q$ are nonempty, compact and decreasing, with diameters $\le2\varepsilon_i\to0$. Their intersection is
a single point $h(p)$, and the characterization holds by construction.

\emph{Continuity.} Fix $p$ and $i$, and let $W$ be the intersection of the links of $D^i$ containing $p$, an open neighbourhood of $p$. For $p'\in W\cap P$ the sets $\st_i(p)$ and
$\st_i(p')$ have a link in common, so $h(p)$ and $h(p')$ lie in the closure of a set of diameter $<4\varepsilon_i$.

\emph{Onto.} Let $q\in Q$ and $i\ge0$, and let $E^i_k\ni q$. As $H_i$ is onto, $k=H_i(j)$ for some $j$, and $D^i_j$ contains a point $p\in P$. Then $q$ and $h(p)$ both lie in
$\cl{\st_i(p)}$, so $\dist(q,h(p))\le2\varepsilon_i$. Hence $h(P)$ is dense in $Q$; it is compact, so $h(P)=Q$.

\emph{One-to-one.} Let $h(p)=h(p')=y$ and $i\ge1$. Since $y\in\cl{\st_i(p)}\cap\cl{\st_i(p')}$ and the closure of a finite union is the union of the closures, $y\in\cl{E^i_a}\cap
\cl{E^i_b}$ with $a=H_i(j)$, $p\in D^i_j$, and $b=H_i(j')$, $p'\in D^i_{j'}$. Since $E^i$ is a BJ-chain, $\abs{a-b}\le1$ (Lemma~\ref{lem:adj-closure}), so $\abs{K_{i-1}(a)-K_{i-1}(b)}\le1$.
By (B), $p\in D^{i-1}_{K_{i-1}(a)}$ and $p'\in D^{i-1}_{K_{i-1}(b)}$. These two links are equal or adjacent, so $\dist(p,p')<2\varepsilon_{i-1}$. As $i$ is arbitrary, $p=p'$.

A continuous bijection of a compact space onto a Hausdorff space is a homeomorphism.
\end{proof}

\formal{\lean{HPC.BJ59.zigzag_homeomorph}. Only the chains $E^i$ need to be BJ-chains.}

\begin{remark}
Bing and Jones's injectivity argument passes through the sets $K_{i-1}\st_i(p)$, whose closures contain $p$ and $q$, and concludes $\rho(p,q)<4/2^i$; with their meshes the bound
should be in terms of $\varepsilon_{i-1}$. In the sentence ``$e_j^{i+1}\subseteq H_iK_i(e_j^{i+1})=H_i(d^{i+1}_k)\subseteq H_id(p,i)$'' of p.~188 the map $H_i$ is applied to a link
of $D^{i+1}$; what is meant is $H_i(K_iH_{i+1}(d^{i+1}_k))$, as in the proof of Lemma~\ref{lem:star-props}(b).
\end{remark}

\section{Fibres}

\begin{proposition}\label{prop:zigzag-fibre}
In the situation of Theorem~\ref{thm:zigzag}, let $\varphi\colon X\to T$ and $\psi\colon Y\to T$ be maps to a set $T$ such that $Q\cap\psi^{-1}(t)$ is closed for every $t$.
Suppose that whenever $p\in P$ lies in $D^i_j$, the link $E^i_{H_i(j)}$ contains a point $q\in Q$ with $\psi(q)=\varphi(p)$. Then $\psi(h(p))=\varphi(p)$ for all $p\in P$.
\end{proposition}

\begin{proof}
For each $i$, $h(p)$ lies in the closure of one of the finitely many links making up $\st_i(p)$, say $E^i_{H_i(j)}$ with $p\in D^i_j$. That link contains $q_i\in Q$ with
$\psi(q_i)=\varphi(p)$, and $\dist(h(p),q_i)\le\varepsilon_i$. So $h(p)$ is a limit of points of the closed set $Q\cap\psi^{-1}(\varphi(p))$.
\end{proof}

\formal{\lean{HPC.BJ59.zigzag_fibre}.}

In Theorem~10, $P=X$, $Q=Y$, $\varphi=f$, $\psi=g$: the hypothesis of Proposition~\ref{prop:zigzag-fibre} is the first half of Condition~A of Chapter~\ref{ch:thm10}, and the
conclusion says $g\circ h=f$; as $h$ is onto, $h(p_x)=q_x$ for every $x$.

\section*{Exercises}

\begin{exercise}
Show that without Condition (B$'$) the sets $\st_i(p)$ need not decrease, and that $h$ may then fail to be well defined.
\end{exercise}

\begin{exercise}
Show that if every $H_i$ and $K_i$ is a bijection preserving the order of the links, then $h$ maps each link-trace $D^i_j\cap P$ into $\cl{E^i_{H_i(j)}}$.
\end{exercise}

\begin{exercise}
Where in the proof of Theorem~\ref{thm:zigzag} is the surjectivity of $H_i$ used? Construct an example with $H_i$ not onto in which $h$ is not onto.
\end{exercise}
